% 
%  chapter7.tex
%  ISELthesis
%  
%  Created by Matilde Pós-de-Mina Pato on 2012/10/09.
%
\chapter{Evaluation}
\label{cha:evaluation}

This chapter evaluates the prototype described in Chapter~\ref{cha:impl} and answers the two
research questions of Section~\ref{sec:int_objectives}: whether internal function calls are bound
correctly at deployment time on both providers (RQ1), and how much that resolution costs (RQ2).
Because the contribution of this work is the \emph{unification} layer (the function registry and
the deployment-time resolution cascade) and the AWS provider added to QuickFaaS, the evaluation
concentrates on the decisions the resolution cascade takes, on the cost of resolving
internal-function endpoints and how that cost scales, and on the size overhead of the AWS agnostic
packaging layer.

\section{Scope and Goals}
\label{sec:eval-scope}

Both the correctness checks and the measurements exercise only \emph{local} operations:
constructing a workflow model, traversing it, resolving internal calls against the registry,
reading and writing the registry file, and packaging a Lambda bundle. Operations that call the
cloud provider (Lambda/IAM/S3 requests, Cloud Run/Cloud Functions APIs, the QuickFaaS subprocess)
are deliberately excluded, for two reasons: they are dominated by network and provider-side latency
measured in seconds, and they are non-deterministic and out of the unit-test scope by design. The
correctness checks replace them with test doubles. For RQ2, the evaluation therefore answers a
focused question: \emph{how much does the unification layer add to a deployment, and how does that
overhead scale with the number of calls and the size of the registry?}

The one exception is analytical rather than measured: \S\ref{sec:eval-api-calls} counts, from the
code, how many provider requests each path through the cascade issues, which is the part of the
cost the local measurements cannot see.

The evaluation has six parts:
\begin{itemize}
  \item \textbf{Correctness of the resolution cascade} (RQ1) --- the decision each production
  resolver takes in every situation the cascade can meet, checked by unit tests on both providers.
  \item \textbf{Resolution overhead and scalability} (RQ2) --- the cost of resolving internal calls as a
  function of the number of calls ($N$) and the size of the registry ($R$), and the effect of the
  single-read optimization.
  \item \textbf{Real auto-deploy resolvers} (RQ2) --- the cost of the production resolvers
  (\texttt{Aws\allowbreak Internal\allowbreak Function\allowbreak Resolver}, \texttt{Workflow\allowbreak Internal\allowbreak Function\allowbreak Resolver}) and of the
  registry write path.
  \item \textbf{Provider API calls} (RQ2) --- how many requests each cascade path sends to the
  provider, counted from the implementation rather than measured.
  \item \textbf{Manual effort} (RQ1) --- the developer actions the integration removes, counted
  against the procedure the two tools required separately.
  \item \textbf{Packaging overhead} --- the size the AWS provider-agnostic layer adds to a Lambda
  deployment bundle.
\end{itemize}

\section{Correctness of the Resolution Cascade}
\label{sec:eval-correctness}

RQ1 asks whether the resolution cascade binds each internal call to the right endpoint and deploys
a function only when its absence is confirmed. Each production resolver
(\texttt{Aws\allowbreak Internal\allowbreak Function\allowbreak Resolver} and
\texttt{Workflow\allowbreak Internal\allowbreak Function\allowbreak Resolver}) has a unit-test
suite that resolves a small workflow in each situation the cascade can meet. The provider is
replaced by test doubles: a function lookup that answers ``found'', ``not found'' or ``access
denied'' per region, a fixed list of regions, and a deployer that records the deployments it is
asked for. The registry is a real file in a temporary directory. Each test checks the endpoint
written into the call, the registry contents afterwards and, where relevant, whether the deployer
was called or the resolution aborted.

Table~\ref{tab:eval-correctness} lists the situations with the behaviour that the cascade
semantics prescribe (\S\ref{subsec:validation-error-semantics}). The observed behaviour matches the
expected behaviour in every case: all 45 resolver tests (22 for AWS, 23 for GCP) pass, as do the
two tests of the registry bootstrap.

\begin{table}[htbp]
  \centering
  \caption{Situations covered by the resolver unit tests, with the expected behaviour. Each row is
  tested on both AWS and GCP unless marked otherwise; every test passes.}
  \label{tab:eval-correctness}
  \begin{tabular}{p{0.42\textwidth}p{0.48\textwidth}}
    \hline
    \textbf{Situation} & \textbf{Expected and observed behaviour} \\
    \hline
    Registry entry valid (Level~1) & call bound to the registered endpoint \\
    Registry entry whose endpoint changed (drift) & registry updated; call bound to the current
    endpoint \\
    Registry entry for a function no longer in its region (stale) & function searched in all
    regions; if found, the entry is replaced; otherwise the entry is removed and the resolution
    aborts \\
    Validation of a registry entry denied by the provider & resolution aborts; entry kept \\
    Not in the registry, found in one region (Level~2) & call bound; function registered \\
    Found in several regions & resolution aborts (ambiguous name) \\
    Explicit \texttt{region/name} reference & only that region is queried \\
    Every candidate region inaccessible, function not found & resolution aborts with a distinct
    error \\
    Not found anywhere, descriptor given (Level~3) & function deployed once from the descriptor;
    registered; call bound \\
    Not found, some regions inaccessible, descriptor given & resolution aborts; nothing deployed \\
    Not found anywhere, no descriptor given & resolution aborts; nothing deployed \\
    Internal calls inside branches, iterations and parallel steps & all resolved (five tests) \\
    External calls & host and path unchanged \\
    Call setting both \texttt{internalFunction} and a literal \texttt{host}/\texttt{path} & rejected
    before the cascade starts \\
    Several internal calls in one workflow & registry read once \\
    First-generation Cloud Function in the registry (GCP only) & call bound without live
    validation (the limitation of \S\ref{subsec:store-internals}) \\
    Registry file missing (provider-independent) & registry rebuilt from the provider's catalog \\
    \hline
  \end{tabular}
\end{table}

The tests check the decisions of the cascade, not the behaviour of the providers: they show that
each resolver acts correctly on what the provider reports, not that the real AWS and GCP APIs
report what the test doubles assume (Section~\ref{sec:eval-threats}).

\section{Methodology}
\label{sec:eval-methodology}

Benchmarks were implemented with the Java Microbenchmark Harness (JMH) 1.37, using
\texttt{AverageTime} mode with results reported in microseconds per operation ($\mu$s/op). Each
measured value is consumed through a \texttt{Blackhole} to prevent dead-code elimination by the
JIT compiler. Every figure reported in this chapter was measured with three JVM forks, three
warm-up iterations and five measurement iterations of one second each
(\texttt{-f 3 -wi 3 -i 5 -w 1 -r 1}), so each value is the mean of fifteen measured iterations
across three JVMs, of already-JIT-compiled code. All experiments but two come from one run of the
whole suite. The exceptions are P10 and P11, which were re-measured later the same evening, on the
same idle machine, after the change to the production resolvers that \S\ref{sec:eval-real-resolvers}
describes; the rest of the suite does not exercise those classes and was left as measured. The two
sessions agree where they can be compared: the $N{=}1$ column of P10 and P11, which is one registry
read and almost nothing else, reproduces the per-read cost that P6 measured in the first session
($9.9$ and $10.4\,\mu$s at $R{=}1$ against P6's $10.5\,\mu$s; $20.0$ and $19.9\,\mu$s at $R{=}50$
against P6's $18.3\,\mu$s).

\paragraph{Caveat on absolute values.}
Measurements were taken on an Apple Silicon laptop running macOS, otherwise idle but not a
dedicated benchmarking machine. With three forks the run-to-run spread is small: the 99.9\%
confidence interval is under 1\% of the score for half the measurements and under 6\% for nine out
of ten. The tables omit those intervals for readability; the few points where the interval is a
large fraction of the score are identified in the text and should be read as noise. What the
figures do not capture is hardware sensitivity: file I/O dominates several of these experiments, so
on a machine with slower storage every registry read and write would cost more, while the
relationships between experiments would hold.

\paragraph{Experiment identifiers.}
Experiments keep the identifiers they carry in the benchmark suite, so the sequence used in this
chapter has gaps. P1, P2, P4, P5 and P12 measure OmniFlow's renderers alone --- rendering time
against the number of functions, the number of call arguments, the AWS and GCP renderers compared,
nesting depth, and the width of a \texttt{Choice} or \texttt{Parallel} --- and therefore exercise
code this work did not change; they are omitted rather than renumbered, so that an identifier here
names the same experiment as in the suite. The chapter reports P3 and P6--P19, which measure
resolution and the registry, plus S1, which measures packaging size.

\paragraph{Notation.}
The experiments share five symbols; what changes between them is the \emph{relationship} among
these quantities, summarised in Table~\ref{tab:eval-notation}.

\begin{table}[htbp]
  \centering
  \caption{Notation used throughout the evaluation.}
  \label{tab:eval-notation}
  \begin{tabular}{cl}
    \hline
    \textbf{Symbol} & \textbf{Meaning} \\
    \hline
    $N$ & total number of \texttt{call} steps in the workflow (internal or external), $N = I + E$ \\
    $R$ & number of entries in the registry file when it is read/written \\
    $F$ & number of \emph{distinct} internal functions the workflow references \\
    $I$ & number of internal calls (repetitions counted) \\
    $E$ & number of external calls (repetitions counted) \\
    \hline
  \end{tabular}
\end{table}

$R$ is the physical size of the registry file and is not the same as $F$: a registry may hold
$R{=}1000$ entries while a workflow references only $F{=}10$ of them. Internal calls are
distributed round-robin over the $F$ distinct functions.

\section{Resolution Overhead and Scalability}
\label{sec:eval-resolution}

\paragraph{Cost of the unification (P3).}
The first experiment isolates the cost of resolving an internal call against the cost of a call
that needs no resolution. It runs the single-read reference resolver over a workflow of $N$ calls
against a one-entry registry ($R{=}1$), once with every call internal and once with every call
external. Both variants pay one registry read, because the resolver reads the file before walking
the tree whether or not anything needs resolving; that read is what the $N{=}1$ figures measure,
about $9.8\,\mu$s (Table~\ref{tab:eval-p3}). What separates the two columns is the per-call work.
The internal path adds a lookup in the in-memory snapshot, the split of the resolved URL into host
and path, and the rebuilding of the call node, which together cost about $0.22\,\mu$s per call,
against about $0.006\,\mu$s for an external call, which is copied unchanged. Resolving a 200-call
workflow therefore costs $54.9\,\mu$s, of which the single read is roughly a fifth. Because the
registry holds one entry, this experiment says nothing about the cost of reading a larger file;
that is the variable the next experiment isolates. The $N{=}10$ point is noise --- its confidence
interval, $\pm16.8\,\mu$s, is wider than the value itself.

\begin{table}[htbp]
  \centering
  \caption{Resolution cost of a workflow of $N$ calls, all internal vs.\ all external, with
  $R{=}1$. Both columns include the one registry read the resolver always performs.}
  \label{tab:eval-p3}
  \begin{tabular}{rrr}
    \hline
    \textbf{Calls ($N$)} & \textbf{Internal --- resolved ($\mu$s)} & \textbf{External --- not resolved ($\mu$s)} \\
    \hline
    1   & 10.1 & 9.8  \\
    2   & 10.4 & 9.8  \\
    5   & 11.1 & 9.8  \\
    10  & 16.9 & 9.8  \\
    20  & 14.5 & 9.8  \\
    50  & 21.3 & 10.0 \\
    100 & 32.2 & 10.3 \\
    200 & 54.9 & 10.9 \\
    \hline
  \end{tabular}
\end{table}

\paragraph{Effect of registry size (P6).}
Fixing the workflow and varying the registry size, with the registry re-read on every call, shows
that the per-call cost grows with $R$ and is essentially independent of $N$: about $10.3\,\mu$s of
fixed I/O and parsing per read plus about $0.17\,\mu$s per registered entry
(Table~\ref{tab:eval-p6}). The four $N$ curves in Figure~\ref{fig:eval-p6} collapse onto a
single line, confirming that the cost grows in direct proportion to both the number of calls and
the registry size. The two terms of that per-call cost are equal when $0.17\,R = 10.3$, that is at
$R \approx 60$: below that size the fixed cost of opening and parsing the file dominates, and only
above it does the number of entries become the larger term. A registry holding the few dozen
functions of a realistic system therefore sits around or below that crossover, where the fixed cost
rules, which is what makes reading it once per workflow --- rather than once per call --- the
change worth making.

\begin{table}[htbp]
  \centering
  \caption{Per-call resolution cost as a function of registry size $R$ (independent of $N$).}
  \label{tab:eval-p6}
  \begin{tabular}{rr}
    \hline
    \textbf{Registry size ($R$)} & \textbf{Cost per call ($\mu$s)} \\
    \hline
    1    & 10.5  \\
    10   & 11.9  \\
    50   & 18.3  \\
    200  & 44.5  \\
    1000 & 183.4 \\
    \hline
  \end{tabular}
\end{table}

\begin{figure}[htbp]
  \centering
  \includegraphics[width=0.8\textwidth]{images/P6_registry_scaling.png}
  \caption{Per-call resolution cost as a function of registry size $R$, for $N \in \{1, 10, 50,
  200\}$ internal calls; the external-call control does not read the registry.}
  \label{fig:eval-p6}
\end{figure}




\paragraph{The single-read optimization (P7--P9).}
Reading the registry once per workflow into an in-memory snapshot and resolving against it
eliminates that call-times-registry-size product: the cost no longer multiplies the two factors
together, becoming instead one read proportional to the registry size plus a cheap constant-time
lookup per call. Table~\ref{tab:eval-p7} compares resolving the registry per call (naive) against a
single read per resolution (optimized). With the snapshot, the cost becomes almost independent of
$N$ --- each row is flat to within a few tenths of a microsecond across a two-hundred-fold increase
in the number of calls --- and reduces to the cost of one read of an $R$-entry file. The speedup
therefore grows with the product $N \cdot R$, reaching $200\times$ in the worst corner
($N{=}200$, $R{=}1000$: $36.7$\,ms down to $0.18$\,ms). Experiments P8 and P9 confirm the same
trend on real \texttt{Workflow} objects rather than the isolated store primitive.

\begin{table}[htbp]
  \centering
  \caption{P7 --- total resolution time, naive (per-call read) vs.\ optimized (single read).}
  \label{tab:eval-p7}
  \begin{tabular}{lrrrr}
    \hline
     & $N{=}1$ & $N{=}10$ & $N{=}50$ & $N{=}200$ \\
    \hline
    naive, $R{=}1$        & 9.8  & 98.0     & 492      & 1\,999    \\
    optimized, $R{=}1$    & 10.0 & 10.0     & 10.1     & 10.3      \\
    naive, $R{=}50$       & 17.9 & 180      & 901      & 3\,630    \\
    optimized, $R{=}50$   & 18.4 & 18.3     & 18.5     & 18.5      \\
    naive, $R{=}1000$     & 183  & 1\,839   & 9\,239   & 36\,670   \\
    optimized, $R{=}1000$ & 182  & 183      & 182      & 184       \\
    \hline
  \end{tabular}
\end{table}

\paragraph{Confirming the optimization on real workflows (P8).}
P8 exercises the resolution of a real \texttt{Workflow} object, not the isolated store primitive,
across a grid of $1$ to $50$ distinct functions and $1$ to $200$ calls, with the registry always
hitting ($R{=}F$). The naive path is dominated by the number of calls, at roughly $10\,\mu$s per
call --- one whole registry read each --- and grows only mildly with $F$, since a larger $F$ means
a larger file to re-read (Table~\ref{tab:eval-p8-naive}). The optimized path pays that read once
and then about $0.22\,\mu$s per call, so it grows with $F$ (the size of the single read) far more
than with $N$ (Table~\ref{tab:eval-p8-optimized}): at $F{=}1$ it goes from $10.2$ to
$54.7\,\mu$s as $N$ goes from 1 to 200, where the naive path goes from $10.4$ to $2\,069\,\mu$s.
The speedup is therefore $1\times$ at $N{=}1$, where both paths read the file exactly once, and
climbs to $38\times$ ($F{=}1$) and $58\times$ ($F{=}50$) at $N{=}200$.

\begin{table}[htbp]
  \centering
  \caption{P8 --- naive resolution time ($\mu$s) of a real workflow, $F$ distinct functions vs.\
  $N$ calls, $R{=}F$.}
  \label{tab:eval-p8-naive}
  \begin{tabular}{rrrrrr}
    \hline
    $F\backslash N$ & 1 & 5 & 10 & 50 & 200 \\
    \hline
    1  & 10.4 & 52.3 & 105 & 525 & 2\,069 \\
    2  & 10.8 & 53.0 & 106 & 534 & 2\,142 \\
    5  & 11.1 & 53.9 & 108 & 543 & 2\,150 \\
    10 & 11.6 & 58.5 & 116 & 582 & 2\,333 \\
    20 & 13.3 & 66.6 & 133 & 665 & 2\,650 \\
    50 & 18.7 & 92.5 & 185 & 929 & 3\,694 \\
    \hline
  \end{tabular}
\end{table}

\begin{table}[htbp]
  \centering
  \caption{P8 --- optimized (single-read) resolution time ($\mu$s) of a real workflow, same grid
  as Table~\ref{tab:eval-p8-naive}.}
  \label{tab:eval-p8-optimized}
  \begin{tabular}{rrrrrr}
    \hline
    $F\backslash N$ & 1 & 5 & 10 & 50 & 200 \\
    \hline
    1  & 10.2 & 11.1 & 12.2 & 20.7 & 54.7 \\
    2  & 10.4 & 11.2 & 12.2 & 21.1 & 54.1 \\
    5  & 10.7 & 11.6 & 12.7 & 21.5 & 54.1 \\
    10 & 11.5 & 12.7 & 13.8 & 22.5 & 55.8 \\
    20 & 13.4 & 14.2 & 15.2 & 24.0 & 57.9 \\
    50 & 18.5 & 19.4 & 20.8 & 29.4 & 63.3 \\
    \hline
  \end{tabular}
\end{table}

On a realistic workflow ($F{=}10$, $N{=}50$ fixed), P9 isolates the effect of the registry size.
Both paths grow with $R$, but for different reasons: the naive path re-reads the file on each of
the 50 calls, the optimized path reads it once, so the gap widens from $26\times$ at $R{=}10$ to
$46\times$ at $R{=}1000$ (Table~\ref{tab:eval-p9}). The optimized column is the cost of one read
plus 50 lookups, and tracks the per-read cost measured in P6.

\begin{table}[htbp]
  \centering
  \caption{P9 --- resolution time ($\mu$s) vs.\ registry size $R$, $F{=}10$/$N{=}50$ fixed.}
  \label{tab:eval-p9}
  \begin{tabular}{rrrr}
    \hline
    \textbf{$R$} & \textbf{naive} & \textbf{optimized} & \textbf{speedup} \\
    \hline
    10   & 589     & 22.6 & 26.1 \\
    20   & 666     & 23.9 & 27.9 \\
    50   & 930     & 29.3 & 31.8 \\
    100  & 1\,378  & 37.8 & 36.5 \\
    200  & 2\,246  & 54.9 & 40.9 \\
    1000 & 9\,487  & 205  & 46.3 \\
    \hline
  \end{tabular}
\end{table}

\section{Real Auto-Deploy Resolvers and Write Path}
\label{sec:eval-real-resolvers}

\paragraph{Production resolvers (P10, P11).}
Experiments P10 and P11 measure the resolvers that actually run during deployment, one per
provider: \texttt{Aws\allowbreak Internal\allowbreak Function\allowbreak Resolver} (Table~\ref{tab:eval-p10}) and
\texttt{Workflow\allowbreak Internal\allowbreak Function\allowbreak Resolver} (Table~\ref{tab:eval-p11}), with the
registry always hitting ($R{=}F$). Both now carry the single-read optimization demonstrated in
P7--P9 (\S\ref{sec:eval-resolution}): each resolver reads the registry once into a snapshot per
\texttt{resolve()} call and looks up every internal call against it
(\texttt{FunctionRegistryStore.tryResolveEntryIn}) instead of re-reading the file per call.
Neither experiment reaches a provider: on AWS the inspector is a local fake that confirms every
registered function, and on GCP the registry holds first-generation \texttt{cloudfunctions.net}
URLs, so the resolver short-circuits before any Cloud Run call
(\S\ref{subsec:store-internals}). What both measure is therefore the local part of the cascade.
The result has the same shape as the already-optimized reference implementation: a cost that is one
registry read plus a small constant per call. Reading down the $N{=}1$ column gives the read, which
grows with $F$ because $R{=}F$; reading along a row gives the per-call work, $0.07$--$0.14\,\mu$s on
AWS and about $0.35\,\mu$s on GCP. Resolving a 200-call workflow against a registry of 10 functions
costs $29.1\,\mu$s on AWS and $81.3\,\mu$s on GCP.

An earlier version of these two experiments measured something else. Until this evaluation was
repeated, both resolvers printed two or three progress lines to the console for every call they
resolved, inside the method being measured, and that console output was about 95\% of the reported
time: the same points cost $1\,006\,\mu$s (AWS) and $1\,344\,\mu$s (GCP). Moving those per-call
messages to \texttt{logger.debug}, which the default logging configuration does not print, produced
the figures above without changing what the resolvers do. The measurement is reported here because
the mistake is easy to repeat: an instrumented method measures its instrumentation, and a
per-call console write costs an order of magnitude more than the resolution it announces.

What remains is a real difference between the providers, not a logging artifact. The AWS resolver
binds a call by concatenating a prefix onto the ARN it already holds; the GCP resolver parses the
resolved URL with \texttt{java.net.URI} and rebuilds host and path from it, once per call. The
reference resolver of P3 and P8 does that same parsing and sits between the two, at
$0.22\,\mu$s per call, which is what places AWS below it and GCP above it.

\begin{table}[htbp]
  \centering
  \caption{P10 --- AWS provider: resolution time ($\mu$s) of the real resolver
  (\texttt{Aws\allowbreak Internal\allowbreak Function\allowbreak Resolver}), $R{=}F$, with the
  single-read optimization and without per-call console output.}
  \label{tab:eval-p10}
  \begin{tabular}{lrrrrr}
    \hline
    $F\backslash N$ & 1 & 5 & 10 & 50 & 200 \\
    \hline
    1  & 9.9  & 10.3 & 10.9 & 13.8 & 24.6 \\
    2  & 10.2 & 11.0 & 11.4 & 14.7 & 26.3 \\
    5  & 11.0 & 11.2 & 11.7 & 15.3 & 27.6 \\
    10 & 11.7 & 12.2 & 12.7 & 16.3 & 29.1 \\
    20 & 14.4 & 14.8 & 14.4 & 19.7 & 34.5 \\
    50 & 20.0 & 20.3 & 20.3 & 26.6 & 47.5 \\
    \hline
  \end{tabular}
\end{table}

\begin{table}[htbp]
  \centering
  \caption{P11 --- GCP provider: resolution time ($\mu$s) of the real resolver
  (\texttt{Workflow\allowbreak Internal\allowbreak Function\allowbreak Resolver}), $R{=}F$, with
  the single-read optimization and without per-call console output.}
  \label{tab:eval-p11}
  \begin{tabular}{lrrrrr}
    \hline
    $F\backslash N$ & 1 & 5 & 10 & 50 & 200 \\
    \hline
    1  & 10.4 & 12.1 & 13.9 & 27.9 & 83.3 \\
    2  & 10.6 & 11.9 & 13.8 & 27.5 & 78.8 \\
    5  & 11.3 & 12.7 & 14.1 & 28.4 & 79.8 \\
    10 & 12.1 & 13.4 & 15.1 & 28.8 & 81.3 \\
    20 & 14.4 & 15.5 & 16.8 & 31.6 & 84.4 \\
    50 & 19.9 & 20.9 & 22.8 & 36.8 & 91.3 \\
    \hline
  \end{tabular}
\end{table}

\paragraph{Registry write path (P13, P17).}
Writing is more expensive than resolving. Every \texttt{put()} re-reads and rewrites the entire
registry file, so registering $K$ functions costs close to $K$ times one write, and one write costs
what it costs to rewrite a file of the current size (Table~\ref{tab:eval-p13}). Both factors are
visible. Across $K$, the cost per write is nearly constant --- at $R_0{=}0$ it moves only from
$45.5$ to $59.1\,\mu$s between $K{=}1$ and $K{=}100$, the mild increase expected from a file that
has gained $K$ entries during the run --- so the total is linear in $K$, not compounding. Across
$R_0$, the cost per write grows roughly in proportion to the file: about $45\,\mu$s on an empty
registry, $114\,\mu$s at $R_0{=}200$ and $490\,\mu$s at $R_0{=}1000$. The registry size is
therefore not a second-order term for writes as it is for reads; it multiplies the cost of every
write.

The miss-then-write pattern the real resolvers perform before any cloud call
(\texttt{tryResolveEntry} followed by \texttt{put}) adds between 26\% and 62\%
(Table~\ref{tab:eval-p17}), because \texttt{tryResolveEntry} reads the whole file once more before
the write rewrites it. The surcharge grows with $R_0$ --- about 30\% on an empty registry, about
60\% at $R_0{=}1000$ --- which is what an extra full read of a growing file predicts. One point
disagrees, $K{=}10$ at $R_0{=}10$, and its confidence interval ($\pm732\,\mu$s on a mean of
$874\,\mu$s) says it is noise rather than an effect. In absolute terms both tables stay in
microseconds per write and reach tens of milliseconds only for $K{=}100$ successive writes into a
thousand-entry registry.

\begin{table}[htbp]
  \centering
  \caption{P13 --- total write time ($\mu$s): $K$ successive \texttt{put} calls, initial registry
  size $R_0$.}
  \label{tab:eval-p13}
  \begin{tabular}{lrrrrr}
    \hline
    $K\backslash R_0$ & 0 & 10 & 50 & 200 & 1000 \\
    \hline
    1   & 45.5     & 48.2     & 63.8     & 114       & 490       \\
    5   & 224      & 242      & 329      & 529       & 1\,945    \\
    10  & 463      & 496      & 640      & 1\,041    & 3\,751    \\
    50  & 2\,517   & 2\,706   & 3\,650   & 5\,363    & 17\,942   \\
    100 & 5\,914   & 6\,693   & 7\,594   & 11\,515   & 36\,195   \\
    \hline
  \end{tabular}
\end{table}

\begin{table}[htbp]
  \centering
  \caption{P17 --- total write time ($\mu$s): $K$ successive miss+\texttt{put} pairs, initial
  registry size $R_0$.}
  \label{tab:eval-p17}
  \begin{tabular}{lrrrrr}
    \hline
    $K\backslash R_0$ & 0 & 10 & 50 & 200 & 1000 \\
    \hline
    1   & 57.2     & 62.7     & 90.2     & 164       & 686       \\
    5   & 298      & 323      & 427      & 769       & 2\,901    \\
    10  & 597      & 874      & 973      & 1\,557    & 5\,751    \\
    50  & 3\,265   & 3\,510   & 5\,051   & 8\,144    & 28\,694   \\
    100 & 8\,369   & 8\,307   & 10\,516  & 17\,689   & 58\,572   \\
    \hline
  \end{tabular}
\end{table}

\section{Match Strategy, Mixed Workflows, and Structure}
\label{sec:eval-secondary}

\paragraph{Exact vs.\ suffix match (P14).}
The suffix-match path (regional disambiguation) fails the exact lookup and scans the loaded map,
reintroducing a cost proportional to the registry size, but over an already-loaded map: the ratio
to exact match grows steadily from $1.28\times$ at $R{=}10$ to $1.95\times$ at $R{=}1000$, and
stays cheap in absolute terms --- $391\,\mu$s for fifty suffix lookups over a thousand-entry
registry (Table~\ref{tab:eval-p14}).

\begin{table}[htbp]
  \centering
  \caption{P14 --- exact vs.\ suffix match ($\mu$s), $F{=}10$/$N{=}50$ fixed.}
  \label{tab:eval-p14}
  \begin{tabular}{rrrr}
    \hline
    \textbf{$R$} & \textbf{exact} & \textbf{suffix} & \textbf{ratio} \\
    \hline
    10   & 22.5 & 28.8 & 1.28 \\
    20   & 24.2 & 32.6 & 1.35 \\
    50   & 29.2 & 40.2 & 1.38 \\
    100  & 38.2 & 58.9 & 1.54 \\
    200  & 55.0 & 92.6 & 1.68 \\
    1000 & 201  & 391  & 1.95 \\
    \hline
  \end{tabular}
\end{table}

\paragraph{Mixed internal/external workflows (P15, P16).}
The first workflows to mix internal and external calls confirm that resolution cost scales with
$I$ (internal calls), not with $N{=}I{+}E$: the $I{=}0$ row stays at $11$--$14\,\mu$s as $E$ grows
to 200, because external calls incur no registry lookup, while reading down a column from $I{=}0$
to $I{=}200$ multiplies the cost by five (Table~\ref{tab:eval-p15}). Each row is flat and each
column grows at about $0.22\,\mu$s per internal call, the same per-call figure as P3. The
registry-size effect proved in P9 is unchanged when a fraction of the calls are external: with
$I{=}40$ internal calls the cost still tracks the one read, from $20\,\mu$s at $R{=}10$ to
$199\,\mu$s at $R{=}1000$ (Table~\ref{tab:eval-p16}).

\begin{table}[htbp]
  \centering
  \caption{P15 --- resolution time ($\mu$s) of a mixed workflow, $R{=}F{=}10$ fixed;
  $I$ internal vs.\ $E$ external calls.}
  \label{tab:eval-p15}
  \begin{tabular}{rrrrrrr}
    \hline
    $I\backslash E$ & 0 & 1 & 5 & 10 & 50 & 200 \\
    \hline
    0   & 11.3 & 11.3 & 11.5 & 11.5 & 13.9 & 12.7 \\
    1   & 11.5 & 11.7 & 11.8 & 11.8 & 14.4 & 12.7 \\
    5   & 12.8 & 12.6 & 12.6 & 12.7 & 12.9 & 13.7 \\
    10  & 13.7 & 13.8 & 13.8 & 13.8 & 14.0 & 14.7 \\
    50  & 22.4 & 22.4 & 22.7 & 26.3 & 22.8 & 23.2 \\
    200 & 55.8 & 55.9 & 56.2 & 67.1 & 57.4 & 56.8 \\
    \hline
  \end{tabular}
\end{table}

\begin{table}[htbp]
  \centering
  \caption{P16 --- resolution time ($\mu$s) vs.\ registry size $R$, mixed workflow with
  $I{=}40$/$E{=}10$/$F{=}10$ fixed.}
  \label{tab:eval-p16}
  \begin{tabular}{rr}
    \hline
    \textbf{$R$} & \textbf{Time ($\mu$s)} \\
    \hline
    10   & 20.4 \\
    20   & 22.2 \\
    50   & 27.3 \\
    100  & 36.2 \\
    200  & 53.7 \\
    1000 & 199  \\
    \hline
  \end{tabular}
\end{table}

\paragraph{Workflow shape (P18, P19).}
Nesting depth (P18) leaves the time essentially constant, $15.8$--$19.2\,\mu$s across depths 0--5
with 20 internal calls held fixed; the depth-2 point ($69.0 \pm 40.8\,\mu$s) is noise, as its own
confidence interval shows. Parallel width (P19) grows only with the total number of internal leaf
calls ($N = 5\times\text{width}$), at $0.22\,\mu$s per call --- the same slope as the flat
workflows of P3 and P15 --- from $12.5\,\mu$s at $N{=}5$ to $119\,\mu$s at $N{=}500$
(Tables~\ref{tab:eval-p18} and~\ref{tab:eval-p19}). Although the resolver recurses explicitly into
\texttt{Iteration} and \texttt{Parallel} contexts, the shape of the tree is irrelevant: only the
count of internal calls matters.

\begin{table}[htbp]
  \centering
  \caption{P18 --- resolution time ($\mu$s) vs.\ nesting depth, 20 internal calls fixed,
  $R{=}F{=}10$.}
  \label{tab:eval-p18}
  \begin{tabular}{rr}
    \hline
    \textbf{Depth} & \textbf{Time ($\mu$s)} \\
    \hline
    0 & 15.9 \\
    1 & 15.8 \\
    2 & 69.0 \\
    3 & 19.2 \\
    4 & 17.1 \\
    5 & 15.9 \\
    \hline
  \end{tabular}
\end{table}

\begin{table}[htbp]
  \centering
  \caption{P19 --- resolution time ($\mu$s) vs.\ \texttt{Parallel} branch width,
  $N = 5\times$ width, $R{=}F{=}10$.}
  \label{tab:eval-p19}
  \begin{tabular}{rrr}
    \hline
    \textbf{Branches} & \textbf{$N$} & \textbf{Time ($\mu$s)} \\
    \hline
    1   & 5   & 12.5 \\
    2   & 10  & 13.7 \\
    5   & 25  & 17.0 \\
    10  & 50  & 22.4 \\
    20  & 100 & 33.2 \\
    50  & 250 & 65.2 \\
    100 & 500 & 119  \\
    \hline
  \end{tabular}
\end{table}

\section{Provider API Calls per Resolution}
\label{sec:eval-api-calls}

The experiments above replace the provider with a local double, so they measure none of the network
traffic the cascade generates. This section counts that traffic rather than measuring it: the
figures are read from the code of the two resolvers and of the clients they use, and they are
numbers of provider API requests, not latencies. Let $G$ be the number of candidate regions --- the
regions enabled for the AWS account, or the Cloud Run locations visible to the GCP project.

\paragraph{What the count grows with.}
Both resolvers validate a registry hit on every internal \texttt{call} step, not once per distinct
function. The snapshot introduced in \S\ref{sec:eval-real-resolvers} removes the repeated
\emph{file} reads, but it does not memoise the live lookup that follows them, so a workflow with
$I$ internal calls that all hit the registry issues $I$ validation requests even when every call
names the same function ($F{=}1$). This is the mechanism behind the growth in $N$, flat in $F$,
that P10 and P11 show for the production resolvers. The list of candidate regions, by contrast, is
fetched at most once per deployment and cached in the resolver, and a region scan is paid only by
the calls that miss.

Table~\ref{tab:eval-api-calls} gives the count for each path through the cascade.

\begin{table}[htbp]
  \centering
  \caption{Provider API requests issued per internal call, by cascade path, counted from the
  implementation. $G$ is the number of candidate regions.}
  \label{tab:eval-api-calls}
  \begin{tabular}{p{0.30\textwidth}p{0.30\textwidth}p{0.30\textwidth}}
    \hline
    \textbf{Path} & \textbf{AWS} & \textbf{GCP} \\
    \hline
    Level~1, registry hit validated & 1 (\texttt{GetFunction} in the region taken from the stored
    ARN) & 1 (Cloud Run \texttt{services.get}) \\
    Level~1, first-generation Cloud Function in the registry & --- & 0 (validation skipped,
    \S\ref{subsec:store-internals}) \\
    Level~1, entry stale & 1 validation $+$ up to $G$ lookups (rediscovery) & 1 validation $+$ up
    to $G$ lookups \\
    Level~2, bare \texttt{functionRef} & up to $G$ \texttt{GetFunction} (the scan stops at the
    second match) & up to $G$ \texttt{services.get} \\
    Level~2, \texttt{"region/name"} reference & 1 & 1 \\
    Level~3, deployment & $G$ lookups (absence is confirmed in every region) $+$ the deployment
    below & $G$ lookups $+$ the deployment below \\
    \hline
  \end{tabular}
\end{table}

Two requests fall outside the per-call count. The candidate-region list costs one EC2
\texttt{DescribeRegions} on AWS, and one Cloud Run \texttt{locations.list} request per response
page on GCP; it is fetched lazily, so a deployment whose calls all resolve at Level~1 never pays
it. The registry bootstrap, which runs only when the registry file is absent, costs that same
region list plus one catalog request per region and per response page --- \texttt{ListFunctions} on
AWS, \texttt{services.list} on GCP.

\paragraph{The deployment itself (Level~3).}
A Level~3 deployment is the only path that writes to the provider, and it is dominated not by the
lookups above but by the deployment sequence. On AWS the resolver resolves the execution role
(\texttt{GetRole}, plus \texttt{CreateRole} and \texttt{AttachRolePolicy} the first time, plus
\texttt{UpdateAssumeRolePolicy} when the trust policy has to be repaired), reads the bucket's
region (\texttt{GetBucketLocation}), runs the QuickFaaS subprocess (which uploads the artifact to
S3 and calls \texttt{CreateFunction}), then polls \texttt{GetFunction} until the function reports
\texttt{Active}, reads the ARN with one further \texttt{GetFunction}, and grants Step Functions
permission to invoke it (\texttt{AddPermission}). On GCP it runs the QuickFaaS subprocess, polls
the Cloud Functions v1 \texttt{get} endpoint until the function is ready, and then calls
\texttt{getIamPolicy} and \texttt{setIamPolicy} to make it invocable. The polls make the total
depend on how long the provider takes to make the function active, which is why this path is
reported as a sequence rather than as a fixed count.

\paragraph{Consequence.}
For a deployment of $I$ internal calls, the traffic ranges from $I$ requests, when every call hits
a valid registry entry, to $1 + I \cdot G$ when the registry is cold and the references are bare
names. The registry is therefore not only a latency optimization but the mechanism that keeps the
request count linear in $I$ instead of proportional to $I \cdot G$; a region-qualified reference
buys the same bound for a call that has to miss.

\section{Manual Effort Before and After}
\label{sec:eval-manual-effort}

The unification is meant to remove manual steps, so it is worth stating which ones, and counting
them. The comparison below is procedural, not a user study: it follows the QuickFaaS procedure
documented in \S\ref{subsec:quickfaas-example} and the integrated procedure of
\S\ref{sec:developer-workflow} for the same task --- making one internal function available to one
workflow on GCP --- and counts the developer actions each requires, and the values the developer
has to carry by hand from one step to the next.

Table~\ref{tab:eval-manual-effort} lists the two sequences. What distinguishes them is not the
total amount of configuration --- the deployment descriptor is written either way --- but the
hand-carried values: in the separate procedure three values produced by one step are read by a
human and typed into the next (the OAuth client ID and secret into the authentication module, the
resulting \texttt{access\_token} into the descriptor, and the generated trigger URL into the
workflow definition), and each is a point where a typo produces a workflow that deploys
successfully and fails at run time. In the integrated procedure that number is zero: the access
token is taken from Application Default Credentials by the deployer, the invoker permission is
granted by the deployer, and the endpoint never leaves the tool --- it goes from the provider into
the registry and from the registry into the rendered workflow.

\begin{table}[htbp]
  \centering
  \caption{Developer actions required to make one function available to one workflow on GCP.
  ``Carried by hand'' counts values produced by one step that the developer must read and re-enter
  in another.}
  \label{tab:eval-manual-effort}
  \begin{tabular}{p{0.45\textwidth}p{0.45\textwidth}}
    \hline
    \textbf{QuickFaaS and OmniFlow used separately} & \textbf{Integrated (this work)} \\
    \hline
    Create an OAuth 2.0 client in the console; copy the client ID and secret \emph{(2 carried by
    hand)} & --- (the deployer uses Application Default Credentials, configured once per machine)
    \\
    Run the auxiliary server, complete the OAuth flow, copy the \texttt{access\_token} into the
    descriptor \emph{(1 carried by hand)} & --- \\
    Write \texttt{func-deployment.json} & Write \texttt{func-deployment.json} (same fields, without
    \texttt{accessToken}) \\
    Write the cloud-agnostic function class & Write the cloud-agnostic function class \\
    Run QuickFaaS to deploy the function & --- (the deployment is triggered by the workflow
    deployment, and only if the function is absent) \\
    Grant the function public or invoker access in the console & --- (granted by the deployer) \\
    Read the generated trigger URL from the console or terminal & --- \\
    Write the workflow, pasting that URL as \texttt{host}/\texttt{path} \emph{(1 carried by hand)} &
    Write the workflow, naming the function:
    \texttt{internalFunction("name", "func-deployment.json")} \\
    Run OmniFlow to deploy the workflow & Run OmniFlow to deploy the workflow \\
    \hline
    \textbf{9 steps, 4 values carried by hand} & \textbf{4 steps, 0 values carried by hand} \\
    \hline
  \end{tabular}
\end{table}

Two qualifications keep the comparison honest. First, the integrated procedure does not remove
configuration, it relocates it: the descriptor is still written by hand, and the credentials still
have to exist --- Application Default Credentials on GCP, environment credentials on AWS --- so
what is saved is the copying, not the setup. Second, the saving is per function and per endpoint
change, so it grows with the number of internal functions a workflow calls and with how often they
are redeployed, and it is largest exactly where the manual procedure is most error-prone: a
workflow that calls several internal functions in several accounts (\S\ref{cha:case-study}).

\section{Packaging Overhead (S1)}
\label{sec:eval-bundle}

The AWS provider added in this work packages a Lambda through the QuickFaaS agnostic layer (the
developer hook plus the \texttt{AwsHttpTemplate} adapter and \texttt{aws-lambda-java-core}). Table
\ref{tab:eval-s1} compares the resulting bundle against an equivalent native \texttt{RequestHandler}
implementation of the same logic. The agnostic layer adds about $0.9$\,KB --- 6.6\% of a trivial
bundle, and well under 0.1\% of a real function --- consistent with the original QuickFaaS
evaluation for GCP and Azure.

\begin{table}[htbp]
  \centering
  \caption{S1 --- Lambda bundle size, AWS agnostic layer vs.\ native implementation.}
  \label{tab:eval-s1}
  \begin{tabular}{lr}
    \hline
    \textbf{Bundle} & \textbf{Size} \\
    \hline
    AWS agnostic (QuickFaaS hook + \texttt{AwsHttpTemplate} + \texttt{aws-lambda-java-core}) & 14\,552 bytes \\
    AWS native (single \texttt{RequestHandler}, same logic) & 13\,647 bytes \\
    \textbf{Delta (agnostic-layer overhead)} & \textbf{905 bytes} \\
    \hline
  \end{tabular}
\end{table}

\section{Discussion}
\label{sec:eval-discussion}

\paragraph{Answer to RQ1.}
In every situation of Table~\ref{tab:eval-correctness}, both resolvers take the expected decision:
each internal call is bound to the registered, current or discovered endpoint, and a function is
deployed only when it is in neither the registry nor any region and every region could be checked.
Stale entries that cannot be rediscovered, ambiguous names and inaccessible regions abort the
resolution instead of binding a call to a guess or deploying a duplicate. The answer to RQ1 is
therefore positive for the resolution logic on both providers, with two qualifications: the
evidence comes from tests that simulate the provider, and on GCP a binding to a first-generation
Cloud Function is trusted without a live check.

\paragraph{Answer to RQ2.}
For RQ2, the benchmarks support six conclusions.

First, the cost of the unification is small per call and linear, and its scaling is well
characterised: a cost that grows in direct proportion to both the number of calls and the
registry size in the naive form, dominated by the fixed cost of reopening the
registry file up to $R \approx 60$ (\S\ref{sec:eval-resolution}). Second, the single-read
optimization eliminates that
call-times-registry-size product, reducing resolution to one read of the registry plus about
$0.22\,\mu$s per internal call, a speedup of $200\times$ in the worst corner measured
(P7, P8, P9) --- a full \emph{identify (P3/P6) $\rightarrow$ fix $\rightarrow$ quantify (P7)} loop.
Third, this optimization has since been propagated to the production resolvers (P10, P11), which
now show the same one-read-plus-a-constant shape as the optimized reference path rather than the
call-times-registry-size cost they used to pay: $29\,\mu$s (AWS) and $81\,\mu$s (GCP) for a
200-call workflow. The difference between the two providers is the per-call URL parsing the GCP
resolver performs and the AWS one does not.
Fourth, the AWS agnostic packaging layer adds negligible size (S1). Fifth, resolution cost tracks
the number of internal calls rather than the workflow's total call count: external calls add only
the shared tree-traversal cost, never a registry lookup (P15, P16). Sixth, neither nesting depth
nor parallel width changes the cost beyond what the total count of internal calls already
predicts: the resolver's explicit recursion into \texttt{Iteration} and \texttt{Parallel} contexts
is free in practice (P18, P19).

\paragraph{Global framing.}
All measured values lie between microseconds and, in the worst write case, tens of milliseconds:
resolving a 200-call workflow costs $55\,\mu$s, and registering a hundred functions in a
thousand-entry registry --- the most expensive operation measured --- costs $36$\,ms. Endpoint
resolution is therefore not the bottleneck of an integrated deployment: the dominant cost is the
actual cloud deployment (seconds). With the single-read
optimization now applied uniformly across both production resolvers, the local overhead of the
unification is asymptotically linear and small beside the deployment it precedes. That conclusion
covers the local work only. The production cascade also validates every registry hit against the
live function and, when a call misses, scans the provider's regions; each of those is a network
round-trip that these benchmarks replace with a local double, so the figures reported here are a
lower bound on what an integrated deployment actually pays for resolution. How many such
round-trips each path costs is counted in \S\ref{sec:eval-api-calls}; how long they take is outside
the scope of this evaluation (\S\ref{sec:eval-threats}).

\section{Threats to Validity}
\label{sec:eval-threats}

The absolute numbers were obtained on a laptop that was idle but not dedicated to benchmarking, so
they characterise this machine rather than serverless deployment hardware in general. The whole
suite was run in one session with three forks, which makes the experiments comparable with one
another and keeps the confidence intervals tight (\S\ref{sec:eval-methodology}); the individual
points whose intervals are wide are identified where they appear and are not used to support any
claim. File I/O dominates several experiments, so storage speed, not processor speed, is the
hardware factor these figures are most sensitive to.

A second limit is one of provenance. P10 and P11 come from a second measurement session, run after
the per-call console output was removed from both resolvers, while every other experiment comes
from the earlier run of the same evening. The two sessions agree on the quantity they share --- the
cost of one registry read (\S\ref{sec:eval-methodology}) --- but the comparison between P10, P11
and the reference path of P8 crosses that boundary and should be read as an order-of-magnitude
comparison rather than an exact one.

The evaluation is also intentionally local-only: it quantifies the resolution and packaging
overhead, not the end-to-end cloud deployment latency, which depends on provider-side behaviour
outside the scope of this work.

The correctness evidence for RQ1 has a similar limit. The resolver tests replace the provider with
test doubles, so they establish that each resolver takes the right decision given what the provider
reports, not that the real AWS and GCP APIs report what the doubles assume. Evidence against the
live providers therefore remains to be produced, and the two full-deployment tests the code base
carries would not supply it: besides needing cloud credentials, which is why they are excluded from
the automated suite, the workflow they deploy gives every call a literal host, so they exercise
workflow deployment and never enter the resolution cascade. Obtaining that evidence means deploying
workflows that contain internal calls against real accounts on both providers, at each level of the
cascade; this is listed as future work (Section~\ref{sec:concl-future}).























































































